% ===== §1 Prelude: the rose in the garden =====
\section{Prelude: The Rose in the Garden}
\label{p09:sec:prelude}

\noindent
In a garden, someone said a small thing that became the seed of this paper. The rose, they remarked, is beautiful; and yet, however hard it labours, it can hold its bloom for only a week or so.

It is the kind of observation one lets pass---a pleasantry about the brevity of beautiful things, the sort of melancholy that attaches to gardens and is half the reason one keeps them. But it lodged, and would not leave, because it names---without meaning to---the precise difficulty this paper is about. The remark seems to set beauty against duration, as though the rose were beautiful \emph{despite} being brief, the brevity a flaw in an otherwise perfect thing. The truth is harder and more interesting: the beauty and the brevity are not two facts but one. \emb{The rose blooms so resplendently \emph{because} it pours the whole of its resource into a single, unrepeatable flowering, holding nothing back against the morrow.} Its resplendence is the very form of its unsustainability. To wish the bloom both as intense and as lasting is to wish a contradiction---to want the flower to spend itself utterly and to keep itself in reserve at once. There is, here, a tragic trade-off that the paper will meet again and again, and that it had better not pretend to abolish: intensity and persistence pull, very often, against one another, and the most intense things are, for that reason, the briefest.

And the same structure, I want to claim, is at work in domains that share no soil with the rose. The intimate bond, at its most incandescent, has the rose's economy: the passion that consumes the whole of itself in its flowering is, precisely thereby, the passion that cannot last, and the lover who demands that the first ardour both blaze and endure demands the rose's contradiction. The logic of capital has it too: the accumulation that drives toward an ever-greater intensity of valorization is the accumulation that exhausts the conditions---the worker, the soil, the future---on which its own continuation depends, blazing brightest as it nears the exhaustion of what feeds it. And the dynamics of the cosmos, at the largest scale, show the same face: the structures that burn most brightly are the most prodigal, and what endures is rarely what shines hardest. Across all three---and the recurrence is itself a reason to attend---one finds a generativity that takes its own continuation as its end, production pointing back at production itself (\zh{反者道之动}, \emph{reversal is the movement of the Dao}), and yet a generativity that may, like the rose, burn itself out in the very intensity by which it lives.

\emc{So the question of this paper is not the one the remark seems to pose.} It is not \emph{how to keep the flower from wilting}---that wish, as we shall see, is the very thing that rots the root. It is the question hidden beneath the remark, and harder to hear: \emph{when the flower has wilted, how does the root go on?} For the observation, true of the flower, is false of the rosebush. The single bloom is mortal and exhaustive, spending all in a week; but the bush flowers year upon year, and it endures not in spite of the flower's death but \emph{by} it---the fallen petals returning their nutrient to the root, this year's wilting the condition of next year's bloom. Here is the distinction on which the whole paper turns, and it is worth stating in its starkest form at the outset:

\begin{claim}[Sustainability is the nesting of finitude, not its abolition]
Sustainability is not the abolition of finitude---not a flower that never falls, which is the bad infinite, the fantasy of perpetual growth that the paper will identify with the vicious circle. Sustainability is the \emph{nesting} of finitude within a higher cycle of regeneration: the flower must fall \emph{so that} the root may go on. A bloom held past its hour does not extend the rose's life; it is the refusal of the fall that kills the bush. The good circle, far from excluding death, requires it---requires that the flower be spent and returned---and what we will come to call the vicious circle is, very often, the one that cannot let its flower fall.
\end{claim}

This reframing carries a consequence the paper will draw out at length, and which it is worth stating now in its sharpest form, for it overturns an assumption so ordinary we scarcely know we hold it. We assume that the threat to a beautiful thing is neglect---that what is precious withers for want of tending, and that more care is therefore always the safer error. The rose teaches otherwise. \emb{The deepest unsustainability comes not from too little tending but, very often, from a tending driven by the fear of unsustainability itself.} The one who, dreading that the rose will wilt, waters it without cease, changes its water, fusses at its roots---that one does not save the flower; he drowns the root, and so manufactures the very withering he feared. This is no incidental cruelty of fate but a structure, and the paper will meet it in every domain it touches: in capital's dread of the end of growth, which drives the extraction that exhausts its base; in the lover's dread of loss, which drives the control and the clutching that rot the bond; in the family's dread of dissolution, which hardens into the very rigidity that breaks it. \zh{为者败之，执者失之}---``the one who acts upon it ruins it, the one who grasps it loses it.'' The fear of an end, made into action, produces the end it fears.

There is, then, a quiet reversal of the imagination at work in this paper from its first page. The one who said the rose blooms but a week was speaking, lovingly, of the flower, and of the sadness of its falling. This paper means to speak of the root---of what is not seen, of what the falling feeds, of the slow regeneration that the flower's brevity makes possible and the anxious hand makes impossible. It is, in a sense, a paper about learning to let the flower fall: not from indifference to its beauty, but from a deeper attention to the root that the beauty, rightly let go, sustains.

\emc{A word, finally, on where this paper stands in the series.} The first eight papers asked, in the main, how a relation \emph{begins}: the mother paper established the ontology of relational being, that the subject is generated in and through relation rather than entering it already complete---a lineage that runs back through Buber's account of the I--Thou as the primary word that constitutes both terms \parencite{buber1970}; \pinkref{Paper~V} descended to the joint attention by which two consciousnesses first constitute one another; \pinkref{Paper~VI} cut into the single decisive moment of the proposal and exposed the epistemic injustice latent within it; \pinkref{Paper~VII} took the dyad's first turning toward the world, the diplomacy by which it faces the external relational field; \pinkref{Paper~VIII} treated the lyric's address to the big Other. Each took, as its object, a \emph{founding moment}---an opening, an inception, a first constitution. \emb{The present paper takes, for the first time, the opposite object: not how a relation begins, but how it does not end.} It turns from inception to maintenance, from the founding act to the long reproduction that must follow it, from the question \emph{how to begin} to the question \emph{how not to end}---which is, we shall find, the question of the root rather than the flower, and the first time the series has asked it. This is its place, and its novelty, in the programme of Generative Relational Being.
